%% ================================================================
%%  SECTION 6 — Judicial Outcome Simulation
%% ================================================================
\section{Judicial Outcome Simulation}

Prediction has the longest research history of any axis in this survey, moving steadily from feature-engineered classifiers through transformer-based encoders to the current crop of legal large language models. The task itself is straightforward to state: take the text of a case, predict whether the appeal was accepted or dismissed, and where possible, generate the reasoning behind that prediction. We are, we believe, at a transition point where the community has begun to recognise that accuracy numbers alone, taken in isolation, mean very little in a legal context. A model that predicts the right outcome for the wrong legal reasons is not a tool; it is a liability.

\subsection{Classical and Early Neural Approaches}

The early period of legal judgment prediction was defined by manual feature engineering. Researchers spent months hand-crafting indicators for bag-of-words and TF-IDF models, and the resulting systems worked, after a fashion. They were also brittle, failing the moment they encountered a case from a different court or a different decade. Neural networks, specifically CNNs and RNNs, removed the need for manual features but ran into a different wall, namely document length. Indian judgments are long. Most early recurrent models either ran out of memory or quietly lost the relevant context by the time they reached the operative order. LegalBERT was the first real breakthrough on this front, demonstrating that domain shift in legal text is severe enough that general-purpose embeddings effectively encode a different language.

\subsection{Transformer-Based Prediction}

Transformers changed what was possible. Through self-attention, models such as \texttt{InLegalBERT}~\cite{paul-etal-2023-inlegalbert}, a BERT-based encoder pre-trained on Indian Supreme Court and High Court judgments, could finally capture dependencies between a factual assertion on page two and a statutory application on page twenty. Foundational benchmarking work by \cite{malik-etal-2021-ildc} established the ILDC corpus, a collection of over thirty-five thousand Supreme Court cases annotated for binary judgment prediction, and made systematic comparison across architectures possible for the first time. The 512-token limit of base BERT, however, remains a persistent frustration in this line of work. We have seen researchers try almost everything (overlapping chunks, sliding windows, hierarchical pooling) but we suspect this ceiling is more damaging than current evaluations reveal — you cannot, in any honest sense, model a hundred-page high court judgment by looking at it through a series of disconnected five-hundred-word keyholes.

\subsection{Explainability as a First-Class Requirement}

A model that outputs ``dismissed'' without an accompanying reason is, for legal use, of little value. You cannot audit it, you cannot identify its biases, and you certainly cannot build trust with it. \texttt{PredEx}~\cite{nigam-etal-2024-predex}, the largest expert-annotated corpus for legal judgment prediction and explanation in the Indian context, addressed this directly by providing roughly fifteen thousand cases in which the rationales are not model-generated hallucinations but actual explanations written by legal experts. This shifts the evaluation question from ``does the output look plausible?'' to ``is the output grounded in expert reasoning?'' We observe, however, that even these expert annotations are often compressed summaries of a much longer reasoning process, and whether a summary can really capture the recursive, self-correcting character of a judge's internal deliberation is, in candour, something we have started to wonder about more often than we used to.

\subsection{Large Language Models and the Scale Frontier}

Scale has become the default answer to almost everything in NLP, and legal AI is no exception. \texttt{NyayaAnumana}, together with its specialised companion model \texttt{INLegalLlama}~\cite{nigam-etal-2025-nyayaanumana}, the largest and most diverse Indian legal corpus compiled to date, currently represents the ceiling for Indian legal judgment prediction, with the language model internalising the rhythms of judicial writing across roughly seven hundred thousand case proceedings spanning multiple court levels. Scale, however, has brought a new problem along with it: benchmark bias. Recent work has shown that many of our current test sets are drawn disproportionately from cases with unambiguous outcomes, which makes models look considerably better in evaluation than they appear to be in the messier reality of a contested courtroom~\cite{nigam-etal-2024-rethinking}. We believe benchmark bias is the single most underappreciated methodological issue in the field today.

\subsection{Remaining Limitations}

Three gaps persist. First, prediction models are still built in isolation from the retrieval and structuring axes discussed in earlier sections. There is no established pathway for a structured fact representation, or a retrieved precedent, to inform a final prediction in any principled way.

Second, the available data is heavily Supreme Court weighted. District courts are where the real backlog lives, yet they have almost no presence in our current benchmarks.

Third, the evaluation metrics are often lazy. Accuracy is the easiest thing to measure, but in a legal context, it is also the most misleading. A model that achieves ninety percent accuracy but fails systematically every time a witness is deemed unreliable is not merely imperfect; it is dangerous, precisely because nobody has thought to check that specific failure mode.

%% ================================================================
%%  SECTION 7 — What-If Analysis
%% ================================================================
\section{What-If Analysis}

What-if analysis is the capability that practitioners most often ask for, and it is also the axis with almost no published literature behind it. It is, in our view, the axis that matters most. Where prediction tells you what the likely outcome is, and explainability tells you why, what-if analysis asks what the outcome would have been had the facts, the statutes invoked, or the precedents relied upon been different. For a lawyer preparing a case, counterfactual reasoning is the form of explainability that actually matters.

\subsection{Counterfactual NLP}

In general NLP, counterfactual analysis usually means perturbing an input to see when the predicted label flips. If a ``dowry demand'' is recharacterised as a ``loan request,'' does the conviction still stand? This is more than a technical curiosity. It is a way to map the decision boundaries of a model and to identify which facts are actually driving the predicted outcome. We distinguish, loosely, between forward-looking what-if analysis (testing hypotheses about how a case might be argued) and backward-looking counterfactuals (explaining past decisions in terms of the minimal changes that would have flipped them).

\subsection{XAI in Legal Decision-Making}

The legal system imposes a transparency requirement that most AI systems are not built to meet. A litigant has, in principle, the right to know why an interest was affected, especially when an algorithm contributed to that decision~\cite{deakin-etal-2023-xai-law}. Most current XAI methods, including LIME and SHAP, are designed to surface feature importance, not the structured counterfactuals that legal reasoning actually requires. Until systems can answer ``what if'' rather than only ``why,'' their explanations will remain descriptive rather than prescriptive.

\subsection{The Indian Gap}

There is no published what-if analysis system for Indian law. We searched, and we found nothing.

This is a substantial opportunity for the community, particularly given the resource constraints typical of Indian litigation, where a tool that lets a lawyer test alternative statutory framings or fact patterns without commissioning a new dataset each time would be genuinely transformative. It is also, we should say, the logical endpoint of the integration we have been arguing for throughout this survey.

Our survey ends here because the literature ends here. We cannot review what has not been built. We can, however, point to the empty space and call it a priority.
